\phantomsection\addcontentsline{toc}{chapter}{课程信息}
\noindent{\bfseries\fontsize{12}{17}\selectfont 课程信息}\par
MIT 18.102/18.1021, Introduction to Functional Analysis, 2021年春季, 本科课程, Casey Rodriguez 授课. 原课堂笔记由 Andrew Lin 整理. 官方\href{https://ocw.mit.edu/courses/18-102-introduction-to-functional-analysis-spring-2021/pages/syllabus/}{大纲}规定每周两次讲课, 每次1.5小时. 页面未公布具体星期、钟点或教室, 本书不补造这些信息.

先修要求为以下两组各一门: 18.06、18.700、18.701 中的一门线性代数课, 以及18.100A、18.100B、18.100P、18.100Q 中的一门实分析课; 也可经教师许可选课. 不指定教材, 官方提供 Richard Melrose 2020年讲义及2021年课堂笔记. 课程由赋范空间、完备性、线性泛函和 Hahn--Banach, 经对偶与有界算子, 进入 Lebesgue 测度、积分、$L^p$ 空间、Hilbert 空间、紧算子和自伴谱理论.

\href{https://ocw.mit.edu/courses/18-102-introduction-to-functional-analysis-spring-2021/pages/calendar/}{日历}列出23讲: 第1--5讲为 Banach 空间基本理论; 第6--12讲为双对偶、外测度、可测集、可测函数和积分收敛; 第13--18讲为 $L^p$、Hilbert、Fourier、投影、Riesz 表示和伴随; 第19--23讲为紧算子、谱和 Dirichlet 问题. 原日历将作业1--5依次列在第2、4、6、8、10讲, 期中列在第11与12讲之间, 作业6--10依次列在第14、16、18、20、22讲, 期末作业列在课程最后. 这是教学次序记录, 不能据此推算未公布的公历截止日期.

作业占总评50\%, 期中25\%, 期末作业25\%. 最低一次作业成绩不计; 评分界限不会高于通常的90分为 A、80分为 B. 作业允许讨论, 但每人须独立写出解答并在右上角注明合作者; 通常不收迟交, 延期须邮件向教师提出. 多数原作业规定经 Gradescope 在24:00前提交; 作业6仅注明 Gradescope, 作业8未写提交说明, 不能据其他份补造其截止时间. 大纲规定第7周有24小时带回完成的期中测验, 课程结束有48小时期末作业.

期中与期末原PDF允许查阅 Richard Melrose 讲义、教师手写笔记、Andrew Lin 笔记、Royden 的 \textit{Real Analysis}（如已购买）、既有作业解答、Piazza 讨论与录播, 禁止与其他学生或互联网协作; 违反者按原文所指 MIT Institute Policy 10.2 处理. 这些是2021年课程的历史考核规则, 本书题解供归档后的学习核对使用.

官方\href{https://ocw.mit.edu/courses/18-102-introduction-to-functional-analysis-spring-2021/pages/assignments-and-exams/}{考核目录}提供10份作业、1份期中和1份 Final Assignment, 共54道编号题、87个显式字母子问及18道无字母单题, 合计105个解答单元. 无字母单题即原编号题未划分(a)、(b)等的题目; 期末第2题虽含两个积分, 按原编号计一个单元, 两个积分均分别解答. 作业10第6题虽注明“不提交”, 仍完整收入. 期末PDF页眉写作 Final, 开头沿用 exam 措辞; 大纲与资源页明确其身份是期末作业, 本书据此命名. 未提供官方答案; 以下新增题解均标为\textbf{AI 独立解答}, 并复用正文已证明的结果. 原例题与习题解答保留, 与官方题目的对应见书末. 本册不以这些解答冒充教师原答案.
